\documentclass[../../main.tex]{subfiles}
\begin{document}

{\bf [解答]}
$n=0$に対しても同様に$f_p(n)$を定義できるので，以下$n\ge0$で考える．
合同式の法を$10$とすると，
\begin{align}
  n^p\equiv f_p(n) \pmod 10
\end{align}
である．故に$k=0,1,2,\cdots,9$および$i\in\mathbb{Z}_{\ge0}$に対して
\begin{align}
  f_p(10i+k)=f_p(k)
\end{align}
が成り立つ．よって以下$f_p(0),\cdots,f_p(9)$についてのみ考えればよい．

\subparagraph{(1)}

$(10-k)^2\equiv k^2$ ゆえ，$f_2(k)=f_2(10-k)$である．故に$k=0,1,2,3,4,5$についてのみ調べ
ればよい．
\begin{align}
  &f_2(0)=0 &f_2(1)=1 \\
  &f_2(2)=4 &f_2(3)=9  \\
  &f_2(4)=6 & f_2(5)=5
\end{align}
だから，求める値は
\begin{align}
  f_2(n)=0,1,4,5,6,9
\end{align}
である．$\cdots$(答)

\subparagraph{(2)}

$n^5-n=n(n-1)(n+1)(n^2+1)$は$2$連続整数の積を因数に含むから$2$の倍数である．

又，ここでのみ合同式の法を$5$とすると，実際に$n\equiv0,\pm1,\pm2,3$を代入すれば
\begin{align}
  n^5 - n \equiv 0 \pmod 5
\end{align}
である．故に$n^5-5$は$5$の倍数である．

以上から，$n^5-n$は$2$かつ$5$の倍数，つまり$10$の倍数である．
よって　
\begin{align}
  n^5-n\equiv 0 \pmod 10 \Longleftrightarrow f_5(n)=f_1(n)
\end{align}
である．$\Box$

\subparagraph{(3)}

(2)の結果から
\begin{align}
  n^{100}\equiv (n^4)^{25}\equiv n^4 \pmod 10
\end{align}
であるから，$f_4{n}$の取り得る値だけを求めればよい．(1)と同様にして，$k=0,1,2,3,4,5$に対して
\begin{align}
  (10-k)^4 \equiv k^4 \pmod 10
\end{align}
だから$k=0,1,2,3,4,5$のみ調べれば十分である．
具体的に計算すると
\begin{align}
  &f_4(0)=0 &f_4(1)=1 \\
  &f_4(2)=6 &f_4(3)=1  \\
  &f_4(4)=6 & f_4(5)=5
\end{align}
であるから，求める値は
\begin{align}
  f_{100}(n)=f_4(n)=0,1,5,6
\end{align}
である．$\cdots$(答)

\newpage
\end{document}
